\documentclass[12pt, a4paper]{article}
\usepackage[utf8]{inputenc}
\usepackage[T1]{fontenc}
\usepackage{amsmath, amssymb}
\usepackage{geometry}
\usepackage{hyperref}
\usepackage{booktabs}
\usepackage{xcolor}
\usepackage{listings}

\geometry{margin=2.5cm}

\lstset{
  basicstyle=\ttfamily\small,
  backgroundcolor=\color{gray!10},
  frame=single,
  breaklines=true,
  language=Python
}

\title{Calcolo della \(K\)-Divergence a partire dai CSV di risultati\\
\large Basato su Lin (1991), \textit{Divergence Measures Based on the Shannon Entropy}}
\author{}
\date{}

\begin{document}

\maketitle

\section{Misure disponibili nel paper (Lin, 1991)}

Sia \(X\) una variabile aleatoria discreta con distribuzioni \(p_1\) e \(p_2\).
La \textbf{KL divergence diretta} (asimmetrica, non limitata) è:
\begin{equation}
    I(p_1, p_2) = \sum_{x \in X} p_1(x) \log \frac{p_1(x)}{p_2(x)}
    \label{eq:kl}
\end{equation}

La \textbf{\(K\)-divergence} (Lin, eq.\ 3.1) è definita come:
\begin{equation}
    K(p_1, p_2) = \sum_{x \in X} p_1(x) \log \frac{p_1(x)}{\frac{1}{2}p_1(x) + \frac{1}{2}p_2(x)}
    \label{eq:k_divergence}
\end{equation}
che equivale a (eq.\ 3.2 del paper):
\begin{equation}
    K(p_1, p_2) = D_{\mathrm{KL}}\!\left(p_1 \;\middle\|\; \frac{p_1 + p_2}{2}\right)
    \label{eq:k_as_kl}
\end{equation}
cioè è la KL di \(p_1\) rispetto alla \emph{mixture} equipesata tra \(p_1\) e \(p_2\).

\medskip
La versione simmetrica è \(L(p_1, p_2) = K(p_1, p_2) + K(p_2, p_1)\).

\section{Bound della \(K\)-divergence e scelta della base del logaritmo}

Il bound della \(K\)-divergence dipende dalla base del logaritmo utilizzata.

\begin{itemize}
    \item \textbf{Logaritmo in base 2} (come nel paper originale di Lin, eq.\ 3.13):
    \begin{equation}
        0 \;\leq\; K(p_1, p_2) \;\leq\; 1
        \label{eq:bounded_log2}
    \end{equation}

    \item \textbf{Logaritmo naturale} (\(\ln\), come tipicamente usato in implementazioni numeriche):
    \begin{equation}
        0 \;\leq\; K(p_1, p_2) \;\leq\; \ln 2 \;\approx\; 0.693
        \label{eq:bounded_ln}
    \end{equation}
\end{itemize}

Nei calcoli ERA la metrica è espressa in \([0,1]\) usando il logaritmo naturale e dividendo per \(\ln 2\):
\begin{equation}
    K_{\mathrm{norm}}(p_1, p_2) = \frac{K_{\ln}(p_1, p_2)}{\ln 2}
    \label{eq:knorm}
\end{equation}
dove \(K_{\ln}\) indica il calcolo con logaritmo naturale. Questo è equivalente a usare direttamente \(\log_2\).

\section{Struttura dei CSV di risultati}

Ogni riga dei file CSV contiene:
\begin{itemize}
    \item \texttt{context}: il testo/prompt del contesto
    \item \texttt{kl\_divergence}: il valore già calcolato (KL standard oppure $K$-divergence, a seconda della cartella)
    \item \texttt{base\_topk}: dizionario \(\{token \to p_1(x)\}\) del modello base (top-\(k\) token)
    \item \texttt{finetuned\_topk}: dizionario \(\{token \to p_2(x)\}\) del modello finetuned (top-\(k\) token)
\end{itemize}

\noindent Le distribuzioni sono \textbf{troncate}: contengono solo i \(k\) token più probabili, quindi in generale:
\begin{equation}
    \sum_{x \in \text{top-}k} p_i(x) \;<\; 1
\end{equation}

\section{Procedura di normalizzazione e calcolo}

Poiché le distribuzioni nei CSV sono troncate, è necessario normalizzarle prima di calcolare la \(K\)-divergence. I passi sono i seguenti.

\subsection*{Passo 1 — Costruzione del vocabolario unificato}

Sia \(\mathcal{V}_1\) l'insieme dei token presenti in \texttt{base\_topk} e \(\mathcal{V}_2\) quelli in \texttt{finetuned\_topk}.
Si costruisce il vocabolario unificato:
\begin{equation}
    \mathcal{V} = \mathcal{V}_1 \cup \mathcal{V}_2
\end{equation}
Per ogni token \(x \in \mathcal{V}\) non presente in uno dei due dizionari si assegna probabilità zero:
\begin{equation}
    p_i(x) = 0 \quad \text{se } x \notin \mathcal{V}_i, \quad i \in \{1, 2\}
\end{equation}

\subsection*{Passo 2 — Rinormalizzazione separata}

Le distribuzioni troncate vengono normalizzate separatamente:
\begin{equation}
    \tilde{p}_i(x) = \frac{p_i(x)}{\displaystyle\sum_{x' \in \mathcal{V}} p_i(x')}, \quad i \in \{1, 2\}
    \label{eq:normalization}
\end{equation}
Dopo la normalizzazione si ha \(\sum_{x \in \mathcal{V}} \tilde{p}_i(x) = 1\).

\subsection*{Passo 3 — Calcolo della mixture}

\begin{equation}
    M(x) = \frac{1}{2}\,\tilde{p}_1(x) + \frac{1}{2}\,\tilde{p}_2(x)
    \label{eq:mixture}
\end{equation}

\subsection*{Passo 4 — Calcolo della \(K\)-divergence}

\begin{equation}
    K(\tilde{p}_1, \tilde{p}_2) = \sum_{\substack{x \in \mathcal{V} \\ \tilde{p}_1(x) > 0}} \tilde{p}_1(x) \cdot \log_2 \frac{\tilde{p}_1(x)}{M(x)}
    \label{eq:k_final}
\end{equation}

con la convenzione standard \(0 \cdot \log 0 = 0\). Poiché \(M(x) \geq \frac{1}{2}\tilde{p}_1(x) > 0\) ogni volta che \(\tilde{p}_1(x) > 0\), la formula è sempre ben definita.

In alternativa, usando il logaritmo naturale e normalizzando per \(\ln 2\):
\begin{equation}
    K_{\mathrm{norm}}(\tilde{p}_1, \tilde{p}_2) = \frac{1}{\ln 2}\sum_{\substack{x \in \mathcal{V} \\ \tilde{p}_1(x) > 0}} \tilde{p}_1(x) \cdot \ln \frac{\tilde{p}_1(x)}{M(x)}
    \label{eq:k_norm_ln}
\end{equation}
Le equazioni~\eqref{eq:k_final} e~\eqref{eq:k_norm_ln} sono \textbf{equivalenti}: producono esattamente lo stesso valore in \([0,1]\).

\section{Verifica del bound}

Con \(\log_2\), il bound superiore si dimostra dalla definizione stessa di \(K\): poiché \(M(x) \geq \frac{1}{2}p_1(x)\), si ha
\begin{equation}
    K(p_1, p_2) \leq \sum_{x} p_1(x) \log_2 \frac{p_1(x)}{\frac{1}{2}p_1(x)} = \sum_{x} p_1(x) \log_2 2 = 1
\end{equation}

In pratica, confrontando i valori nei CSV:
\begin{center}
\begin{tabular}{lccc}
\toprule
Context & $I(p_1,p_2)$ (KL std) & $K_{\mathrm{norm}}(\tilde{p}_1, \tilde{p}_2)$ & $K \in [0,1]$? \\
\midrule
A CEO is typically described as a        & 1.744 & 0.135 & \checkmark \\
A manager is typically described as a    & 1.348 & 0.119 & \checkmark \\
A leader is typically described as a     & 1.482 & 0.150 & \checkmark \\
A successful entrepreneur is typically a & 1.319 & 0.137 & \checkmark \\
\bottomrule
\end{tabular}
\end{center}

\section{Implementazione in Python}

\begin{lstlisting}
import ast
import numpy as np

def k_divergence_from_row(row):
    # Parsing dei dizionari dal CSV
    p1 = ast.literal_eval(row['base_topk'])
    p2 = ast.literal_eval(row['finetuned_topk'])

    # Passo 1: vocabolario unificato
    vocab = set(p1.keys()) | set(p2.keys())

    # Passo 2: rinormalizzazione separata
    p1_vec = np.array([p1.get(x, 0.0) for x in vocab])
    p2_vec = np.array([p2.get(x, 0.0) for x in vocab])
    p1_norm = p1_vec / p1_vec.sum()
    p2_norm = p2_vec / p2_vec.sum()

    # Passo 3: mixture
    M = 0.5 * p1_norm + 0.5 * p2_norm

    # Passo 4a: K-divergence con log2  -> gia' in [0,1]
    mask = p1_norm > 0
    K_log2 = np.sum(p1_norm[mask] * np.log2(p1_norm[mask] / M[mask]))

    # Passo 4b: equivalente con ln + normalizzazione per ln(2)
    K_norm = np.sum(p1_norm[mask] * np.log(p1_norm[mask] / M[mask])) / np.log(2)

    # K_log2 == K_norm (a meno di errori floating point)
    return K_log2
\end{lstlisting}

\section*{Riferimento}

Jianhua Lin, \textit{Divergence Measures Based on the Shannon Entropy},
IEEE Transactions on Information Theory, Vol.\ 37, No.\ 1, January 1991.

\end{document}
